\section{第6讲\quad 基础测试卷}

本试卷共9个题目，选择5个，填空3个，解答1个，均为基础水平题目。限时30分钟。得分在70分以上视为合格，可以去做【刷中档】模块题目，冲刺更高的水平。

\subsection{一、单选题，共5题，每题10分}

\begin{example}[\bkstar{1}{5}]
下列命题中正确的有

①一个棱柱至少有 $5$ 个平面；

②正棱锥的侧面都是全等的等腰三角形；

③有两个面平行且相似，其他各面都是梯形的多面体是棱台；

④正方形的直观图是正方形．
\begin{tasks}(4)
    \task $1$
    \task $2$
    \task $3$
    \task $4$
\end{tasks}
\end{example}
\begin{solution}
\textbf{答案：}B

①因为底面最少为三角形，故 $3$ 个侧面，$2$ 个底面，共 $5$ 个面，故①对

②正棱锥的底面是正多边形，顶点在底面上的射影是底面正多边形的中心，射影侧面都是全等的等腰三角形，②正确

③不正确，因为不能保证各侧棱的延长线交与一点

④正方形的直观图是平行四边形

所以④不正确

正确的命题只有①②

故选：B．
\end{solution}

\vspace{2em}

\begin{example}[\bkstar{1}{5}]
已知某圆柱的正视图是面积为 $4$ 的正方形，则此圆柱的体积为
\begin{tasks}(4)
    \task $\pi$
    \task $2\pi$
    \task $3\pi$
    \task $4\pi$
\end{tasks}
\end{example}
\begin{solution}
\textbf{答案：}B

$\because$ 圆柱的正视图是面积为 $4$ 的正方形

$\therefore$ 圆柱的底面半径为 $1$，高为 $2$

$\therefore$ 圆柱的体积 $V=\pi\times 1^2\times 2=2\pi$

故选：B．
\end{solution}

\vspace{2em}

\begin{example}[\bkstar{1}{5}]
已知一个球的体积为 $\dfrac{4}{3}\pi$，则该球的表面积为
\begin{tasks}(4)
    \task $\pi$
    \task $2\pi$
    \task $3\pi$
    \task $4\pi$
\end{tasks}
\end{example}
\begin{solution}
\textbf{答案：}D

一个球的体积为 $\dfrac{4}{3}\pi$，则 $\dfrac{4}{3}\pi r^3=\dfrac{4}{3}\pi$，解得半径 $r=1$

得表面积为：$4\pi r^2=4\pi$

故选：D．
\end{solution}

\vspace{2em}

\begin{example}[\bkstar{1}{5}]
下列各个条件中，可以确定一个平面的是
\begin{tasks}(2)
    \task 三个点
    \task 两条不重合直线
    \task 一个点和一条直线
    \task 不共点的两两相交的三条直线
\end{tasks}
\end{example}
\begin{solution}
\textbf{答案：}D

A．不共线的三点确定一个平面，如果三点共线，不能确定一个平面，故A错误

B．两条直线异面时，不能确定一个平面，故B错误

C．若点在直线上，则不能确定一个平面，故C错误

故选：D．
\end{solution}

\vspace{2em}
\begin{example}[\bkstar{2}{5}]
经过平面外两点可作与这个平面平行的平面的个数为
\begin{tasks}(4)
    \task $0$ 个
    \task $1$ 个
    \task $0$ 或 $1$ 个
    \task $1$ 或 $2$ 个
\end{tasks}
\end{example}
\begin{solution}
\textbf{答案：}C

若这两点的连线与该平面相交，则不能作出符合条件的平面

若这两点的连线与该平面平行，则可作出一个符合条件的平面

故选：C．
\end{solution}

\vspace{2em}

\subsection{二、填空题，共3题，每题10分}

\begin{example}[\bkstar{2}{5}]
已知三个不同的平面 $\alpha$，$\beta$，$\gamma$ 满足 $\alpha\perp\gamma$，$\beta\perp\gamma$，则 $\alpha$ 与 $\beta$ 的关系是
\end{example}
\begin{solution}
\textbf{答案：}相交或平行

如图
\bkfig{TikZ-final/geometry/solid/p096-o1842.tex}

在正方体 $ABCD-A_1B_1C_1D_1$ 中，平面 $ADD_1A_1\perp$ 平面 $ABCD$，平面 $DCC_1D_1\perp$ 平面 $ABCD$

平面 $ADD_1A_1\cap$ 平面 $DCC_1D_1=DD_1$

平面 $ADD_1A_1\perp$ 平面 $ABCD$，平面 $BCC_1B_1\perp$ 平面 $ABCD$，平面 $ADD_1A_1\px$ 平面 $BCC_1B_1$

所以三个不同的平面 $\alpha$，$\beta$，$\gamma$ 满足 $\alpha\perp\gamma$，$\beta\perp\gamma$，$\alpha$ 与 $\beta$ 相交或平行

故答案为：相交或平行
\end{solution}

\vspace{2em}

\begin{example}[\bkstar{2}{5}]
已知正三棱锥 $S-ABC$ 的侧棱长为 $2$，底面边长为 $1$，则侧棱 $SA$ 与底面 $ABC$ 所成角的余弦值等于
\end{example}
\begin{solution}
\textbf{答案：}$\dfrac{\sqrt{3}}{6}$

如图所示，作 $SO\perp$ 平面 $ABC$，则 $O$ 是底面等边三角形的中心
\bkfig{TikZ-final/geometry/solid/p096-o1847.tex}

故 $\angle SAO$ 是直线 $SA$ 与底面 $ABC$ 所成的角，其中 $AO=\dfrac{\sqrt{3}}{3}AB=\dfrac{\sqrt{3}}{3}$，$SA=2$

故 $\cos\angle SAO=\dfrac{AO}{SA}=\dfrac{\frac{\sqrt{3}}{3}}{2}=\dfrac{\sqrt{3}}{6}$

故答案为：$\dfrac{\sqrt{3}}{6}$．
\end{solution}

\vspace{2em}

\begin{example}[\bkstar{3}{5}]
已知正方体 $ABCD-A_1B_1C_1D_1$ 的棱长为 $2$，则点 $D$ 到平面 $ACD_1$ 的距离为
\begin{tasks}(4)
    \task $\dfrac{2\sqrt{3}}{9}$
    \task $\dfrac{2\sqrt{3}}{3}$
    \task $\sqrt{3}$
    \task $\dfrac{4\sqrt{3}}{3}$
\end{tasks}
\end{example}
\begin{solution}
\textbf{答案：}B

如图，依题意知 $DD_1\perp$ 平面 $ADC$
\bkfig{TikZ-final/geometry/solid/p097-o1860.tex}

则 $V_{D_1-ADC}=\dfrac{1}{3}\times\dfrac{1}{2}\times 2\times 2\times 2=\dfrac{4}{3}$

$\because AD_1=AC=CD_1=2\sqrt{2}$

$\therefore S_{\triangle ACD_1}=\dfrac{\sqrt{3}}{4}\times\left(2\sqrt{2}\right)^2=2\sqrt{3}$

设 $D$ 到平面 $ACD_1$ 的距离为 $d$

则 $V_{D-ACD_1}=\dfrac{1}{3}\cdot d\cdot S_{\triangle ACD_1}=\dfrac{1}{3}\cdot d\cdot 2\sqrt{3}$

由 $V_{D_1-ADC}=V_{D-AD_1C}$，可得 $d=\dfrac{2\sqrt{3}}{3}$

故选：B．
\end{solution}

\vspace{2em}
\subsection{三、解答题，共1题，每题20分}

\begin{example}[\bkstar{2}{5}]
如图，在直三棱柱 $ABC-A_1B_1C_1$ 中，$BC\perp AC$，$D$，$E$ 分别是 $AB$，$AC$ 的中点．
\begin{enumerate}
    \item[(1)] 求证：$B_1C_1\px$ 平面 $A_1DE$；
    \item[(2)] 求证：平面 $A_1DE\perp$ 平面 $ACC_1A_1$．
\end{enumerate}
\bkfig[0.3\textwidth]{TikZ-final/geometry/solid/p042-o0861.tex}
\end{example}
\begin{solution}
(1) 因为 $D$，$E$ 分别是 $AB$，$AC$ 的中点

所以 $DE\px BC$

又因为在三棱柱 $ABC-A_1B_1C_1$ 中，$B_1C_1\px BC$

所以 $B_1C_1\px DE$

又 $B_1C_1\not\subset$ 平面 $A_1DE$，$DE\subset$ 平面 $A_1DE$

所以 $B_1C_1\px$ 平面 $A_1DE$

(2) 在直三棱柱 $ABC-A_1B_1C_1$ 中，$CC_1\perp$ 底面 $ABC$

又 $DE\subset$ 底面 $ABC$

所以 $CC_1\perp DE$

又 $BC\perp AC$，$DE\px BC$

所以 $DE\perp AC$

又 $CC_1$，$AC\subset$ 平面 $ACC_1A_1$，且 $CC_1\cap AC=C$

所以 $DE\perp$ 平面 $ACC_1A_1$

又 $DE\subset$ 平面 $A_1DE$

所以平面 $A_1DE\perp$ 平面 $ACC_1A_1$
\end{solution}

\vspace{2em}
